\documentclass[10pt,a4paper]{amsart}
\usepackage[margin=19mm]{geometry}
\usepackage{amsmath,amssymb,mathtools}
\usepackage[scr=rsfs,cal=euler]{mathalfa}
\usepackage{xfrac}
\usepackage[unicode,hidelinks,bookmarks=false]{hyperref}
\newcommand{\ie}{i.e.\ }
\newcommand{\cf}{cf.\ }
\newcommand{\resp}{respectively\ }
\newcommand{\Subsection}{\subsection}
\providecommand{\nobreakdash}{\nobreak\hbox{-}}
\setlength{\emergencystretch}{2em}
\multlinegap0pt
\usepackage{amssymb}
\usepackage{mathtools}
\newtheorem{lemma}{Lemma}
\newcommand{\Z}{\mathbb Z}
\title{Prime-power Diophantine tuples}
\author{Andrej Dujella}
\date{Source: arXiv:2609.03448v1 (CC BY 4.0)}
\begin{document}
\maketitle

\noindent\emph{Source and licence.} Prime-power Diophantine tuples. Version arXiv:2609.03448v1 (CC BY 4.0), distributed by arXiv under the Creative Commons Attribution 4.0 International licence, http://creativecommons.org/licenses/by/4.0/, as recorded in the arXiv licence metadata for this exact version and shown on the abstract page https://arxiv.org/abs/2609.03448v1. Full licence notice: \texttt{licenses/arxiv-2609.03448-CC-BY-4.0.txt}.\par\smallskip

\noindent\textit{Editorial scope.} Counted unit: the article's lemma lem:toric (source0.tex lines 510--541). The article's complete original proof of it is delivered with it (source0.tex lines 543--626, one \texttt{proof} environment of 84 lines). No mathematics is written, repaired or re-derived: the statement and the proof are the article's own lines, in the article's own notation, and no retained line is substituted or re-flowed.  Retained uncounted as the source-local context the proof needs: lines 497--499.  Nothing was omitted from any retained span, so the package carries no omission record.  No retained line contains a citation, so the wrapper carries no bibliography and no bibliography entry.  No retained line contains a figure environment, an image-inclusion command or drawing code, so no image is delivered and the package declares no asset dependency.  Editorial formatting only, in addition: an AMS \texttt{amsart} preamble that restates the article's own declarations and macros used by the retained lines, namely the environments \texttt{lemma} on separate counters, together with the standard packages those lines need.  The retained environments print in this document as Lemma 1 in order of appearance, whereas the article prints the same items with the numbering of its own full document.  The complete native source of the article as distributed in the arXiv e-print for this exact version is bundled unchanged as \texttt{sources/209297/source0.tex}.
\begin{equation}\label{eq:AiBi}
 A_i(t)\equiv T_iB_i(t)\pmod N,
\end{equation}

\begin{lemma}[toric eliminant]\label{lem:toric}
Let
\[
 D_0=24^4=331776.
\]
There exists a nonconstant polynomial
\[
 Q(T_1,\ldots,T_5)\in\Z[T_1,\ldots,T_5]
\]
with
\begin{equation}\label{eq:Q-degree}
 1\le q:=\deg Q\le5D_0<2^{21},
\end{equation}
and polynomials
\[
 G_i\in\Z[t_1,\ldots,t_4,T_1,\ldots,T_5]\qquad(1\le i\le5)
\]
such that
\begin{equation}\label{eq:toric-certificate}
 (t_1t_2t_3t_4)^{12D_0}Q(T)
 =\sum_{i=1}^5G_i(t,T)\bigl(A_i(t)-T_iB_i(t)\bigr).
\end{equation}
The same polynomial $Q$ may be chosen so that
\begin{equation}\label{eq:Q-height-combined}
 \log_2 \bigl(h(Q)64^q\bigr)<2^{116}.
\end{equation}
Consequently every reduced labelled quintuple satisfying~\eqref{eq:AiBi}
satisfies
\begin{equation}\label{eq:Q-cong}
 Q(64\Lambda_1,\ldots,64\Lambda_5)\equiv0\pmod N.
\end{equation}
\end{lemma}

\begin{proof}
For $\alpha=(\alpha_1,\ldots,\alpha_5)$ with
$0\le\alpha_i\le D_0$, define
\[
 F_\alpha(t)=\prod_{i=1}^5A_i(t)^{\alpha_i}B_i(t)^{D_0-\alpha_i}.
\]
For a fixed variable $t_j$, only the four factors indexed by $i\ne j$
involve $t_j$. In each such factor, the $t_j$-degree is at most
\[
 6\alpha_i+3(D_0-\alpha_i)\le6D_0.
\]
Hence
\[
 \deg_{t_j}F_\alpha\le24D_0\qquad(j=1,\ldots,4).
\]
Therefore, the coefficient vectors of all $F_\alpha$ lie in a vector space of
dimension at most
\[
 M=(24D_0+1)^4.
\]
There are $(D_0+1)^5$ such polynomials. Since $D_0=24^4$,
\[
 (24D_0+1)^4<[24(D_0+1)]^4
 =D_0(D_0+1)^4<(D_0+1)^5.
\]
Thus the family is linearly dependent.

Let $\rho$ be the rank of its integer coefficient matrix. Choose $\rho$
linearly independent columns and one further column in their span. Choose a
nonzero $\rho\times\rho$ minor of the independent columns and use the
corresponding signed maximal minors to obtain an integral relation
\begin{equation}\label{eq:F-relation}
 \sum_\alpha c_\alpha F_\alpha(t)=0
\end{equation}
supported on at most $\rho+1$ columns. Using these same coefficients, define
\[
 Q(T)=\sum_\alpha c_\alpha T_1^{\alpha_1}\cdots T_5^{\alpha_5}.
\]
The monomials $T^\alpha$ are distinct, so $Q\ne0$. Moreover $Q$ is
nonconstant: a constant relation would be supported only at
$\alpha=(0,\ldots,0)$, whereas
\[
 F_{\mathbf0}=\prod_{i=1}^5B_i^{D_0}\ne0.
\]
This proves~\eqref{eq:Q-degree}.

Put $S(t)=\prod_iB_i(t)^{D_0}$. Since
\[
 \prod_{i=1}^5B_i(t)=(t_1t_2t_3t_4)^{12},
\]
we have $S=(t_1t_2t_3t_4)^{12D_0}$. For each $\alpha$,
$F_\alpha-ST^\alpha$ belongs to the ideal generated by the five polynomials
$A_i-T_iB_i$. Indeed, this follows by replacing the factors one at a time and
using $U^m-V^m=(U-V)(U^{m-1}+\cdots+V^{m-1})$. Summing with the coefficients
$c_\alpha$ in~\eqref{eq:F-relation} gives~\eqref{eq:toric-certificate}.
Since $t_1t_2t_3t_4$ is a unit modulo $N$,~\eqref{eq:Q-cong} follows.

It remains to bound the same maximal-minor relation. Each $\Delta_i$ is a
product of six binomials, so
\[
 h(A_i)\le2^{12},\qquad h(B_i)=1,
\]
and therefore $h(F_\alpha)\le2^{60D_0}$. The Euclidean norm of every
coefficient column is at most its $\ell^1$-norm, hence at most $2^{60D_0}$.
Hadamard's inequality gives
\[
 h(Q)\le(\rho+1)2^{60D_0\rho}.
\]
Furthermore,
\[
 \rho\le(24D_0+1)^4<2^{92},\qquad q\le5D_0.
\]
Hence,
\[
 \begin{split}
 \log_2 \bigl(h(Q)64^q\bigr)
 &\le \log_2(\rho+1)+60D_0\rho+6q\\
 &<92+60D_0(24D_0+1)^4+30D_0\\
 &<2^{116},
 \end{split}
\]
where the last inequality is a direct integer calculation. This proves
\eqref{eq:Q-height-combined}.
\end{proof}

\end{document}
